\chapter{Concentration of Measures}
\section{Lecture 20: Generalization and Transformations}


We first motivate why we care about the concentration inequalities:

The main themes are:
\begin{enumerate}[label=(\arabic*)]
    \item concentration bounds imply consistency and stochastic rates;
    \item sub-Gaussian and sub-exponential classes encode useful tail behavior;
    \item Bernstein-type inequalities control sums and sample means;
    \item bounded differences give concentration for functions of independent samples.
\end{enumerate}

\begin{prop}[Consistency and stochastic rates]
Let $\hat\theta_n$ be an estimator of $\theta$.

\begin{enumerate}[label=(\alph*)]
    \item Suppose that for every $t>0$,
    \[
    \pr\bigl(|\hat\theta_n-\theta|>t\bigr)\le a_n(t),
    \qquad a_n(t)\to 0 \text{ as } n\to\infty.
    \]
    Then $\hat\theta_n \to \theta$ in probability.

    \item Suppose further that there exists a deterministic sequence $r_n\to\infty$ such that for every $M>0$,
    \[
    \pr\bigl(r_n|\hat\theta_n-\theta|>M\bigr)
    \le a_n(M/r_n)\to 0.
    \]
    Then
    \[
    r_n(\hat\theta_n-\theta)=O_p(1).
    \]
\end{enumerate}
\end{prop}

\begin{proof}
For (a), fix $t>0$. By assumption,
\[
\pr\bigl(|\hat\theta_n-\theta|>t\bigr)\le a_n(t)\to 0,
\]
which is exactly convergence in probability.

For (b), fix $M>0$. Then
\[
\pr\bigl(|r_n(\hat\theta_n-\theta)|>M\bigr)
=
\pr\bigl(|\hat\theta_n-\theta|>M/r_n\bigr)
\le a_n(M/r_n).
\]
By assumption, this tends to $0$ for each fixed $M$, which is the definition of $O_p(1)$.
\end{proof}

\begin{remark}
A concentration inequality is a \emph{non-asymptotic} form of error control.  
Consistency says the estimation error vanishes; a sharper concentration bound says \emph{how fast} and with \emph{what tail behavior}.
\end{remark}



\begin{prop}[Equivalent characterizations of sub-Gaussianity]
For a centered random variable $X$, the following are equivalent up to universal constants:
\begin{enumerate}[label=(\roman*)]
    \item There exists $K_1>0$ such that
    \[
    \E e^{\lambda X}\le \exp\!\left(\frac{K_1^2\lambda^2}{2}\right)
    \qquad \forall \lambda\in\R.
    \]
    \item There exists $K_2>0$ such that
    \[
    \pr(|X|>t)\le 2e^{-t^2/K_2^2}
    \qquad \forall t\ge 0.
    \]
    \item There exists $K_3>0$ such that
    \[
    (\E |X|^p)^{1/p}\le K_3\sqrt{p}
    \qquad \forall p\ge 1.
    \]
    \item There exists $K_4>0$ such that
    \[
    \E \exp(X^2/K_4^2)\le 2.
    \]
\end{enumerate}
\end{prop}
\begin{proof}
Assume throughout that $\E X=0$ when discussing condition \emph{(i)}.
If $X$ is not centred, then one should replace $X$ by $X-\E X$ in \emph{(i)}.

\medskip
\noindent
\textbf{(i)$\Rightarrow$(ii).}
Fix $t>0$ and $\lambda>0$. By Markov's inequality and the condition \emph{(i)},
\[
\pr(X>t)
=
\pr\bigl(e^{\lambda X}>e^{\lambda t}\bigr)
\le
e^{-\lambda t}\E e^{\lambda X}\le\exp\!\left(-\lambda t+\frac{K_1^2\lambda^2}{2}\right).
\]
The right-hand side is minimized at
\[
\lambda=\frac{t}{K_1^2},
\]
and substituting this value gives
\[
\pr(X>t)\le \exp\!\left(-\frac{t^2}{2K_1^2}\right).
\]
Applying the same argument to $-X$ and the union bound gives
\[
\pr(|X|>t)
\le
2\exp\!\left(-\frac{t^2}{2K_1^2}\right).
\]
Thus \emph{(ii)} holds with $K_2=\sqrt{2}\,K_1$.

\medskip
\noindent
\textbf{(ii)$\Rightarrow$(iii).}
Let $p\ge 1$. By Tonelli's theorem,
\begin{align*}
\mathbb{E}|X|^p&=\mathbb{E} \int_0^{\infty} \mathbf{1}_{\{s<|X|^p\}} d s=\int_0^{\infty} \mathbb{P}(|X|^p>s) d s\\
&=\int_0^{\infty} \mathbb{P}\left(|X|>s^{1 / p}\right) d s
\end{align*}
Now make the change of variables
\begin{align*}
s=t^p, \quad d s=p t^{p-1} d t
\end{align*}
we then have
\begin{align*}
\mathbb{E}|X|^p=\int_0^{\infty} p t^{p-1} \mathbb{P}(|X|>t) d t
\end{align*}
By \emph{(ii)},
\[
\E |X|^p
\le
2p\int_0^\infty t^{p-1}e^{-t^2/K_2^2}\,dt.
\]
Make the change of variable
\[
u=\frac{t^2}{K_2^2},
\qquad
t=K_2\sqrt{u},
\qquad
dt=\frac{K_2}{2}u^{-1/2}\,du.
\]
Then
\begin{align*}
\E |X|^p
&\le
2p\int_0^\infty (K_2\sqrt{u})^{p-1}e^{-u}\frac{K_2}{2}u^{-1/2}\,du \\
&=
pK_2^p\int_0^\infty u^{p/2-1}e^{-u}\,du \\
&=
pK_2^p\,\Gamma\!\left(\frac p2\right)
=
2K_2^p\,\Gamma\!\left(\frac p2+1\right).
\end{align*}
Now use the standard bound
\[
\Gamma(r+1)\le (Cr)^r,
\qquad r\ge \frac12,
\]
for some universal constant $C>0$. With $r=p/2$ this gives
\[
\E |X|^p
\le
2K_2^p \left(C\frac p2\right)^{p/2}
\le
(C'K_2\sqrt p)^p
\]
for another universal constant $C'$. Taking $p$th roots,
\[
(\E |X|^p)^{1/p}\le C'K_2\sqrt p.
\]
Thus \emph{(iii)} holds with $K_3\le C'K_2$.

\medskip
\noindent
\textbf{(iii)$\Rightarrow$(iv).}
Expand the exponential into its power series:
\[
\E e^{X^2/L^2}
=
\sum_{m=0}^\infty \frac{\E X^{2m}}{m!\,L^{2m}}.
\]
By \emph{(iii)}, with $p=2m$,
\[
(\E |X|^{2m})^{1/(2m)}\le K_3\sqrt{2m},
\]
hence
\[
\E X^{2m}\le (K_3\sqrt{2m})^{2m}=K_3^{2m}(2m)^m.
\]
Therefore
\[
\frac{\E X^{2m}}{m!\,L^{2m}}
\le
\frac{K_3^{2m}(2m)^m}{m!\,L^{2m}}.
\]
Choose
\[
L=2\sqrt e\,K_3.
\]
Then
\[
L^{2m}=(4eK_3^2)^m,
\]
so
\[
\frac{\E X^{2m}}{m!\,L^{2m}}
\le
\frac{(2m)^m}{m!(4e)^m}
=
\frac{1}{m!}\left(\frac{m}{2e}\right)^m.
\]
Using Stirling's lower bound
\[
m!\ge \left(\frac me\right)^m,
\]
we obtain
\[
\frac{\E X^{2m}}{m!\,L^{2m}}
\le 2^{-m}.
\]
Hence
\[
\E e^{X^2/L^2}
\le
\sum_{m=0}^\infty 2^{-m}
=2.
\]
So \emph{(iv)} holds with
\[
K_4\le 2\sqrt e\,K_3.
\]

\medskip
\noindent
\textbf{(iv)$\Rightarrow$(ii).}
For any $t\ge 0$, Markov's inequality gives
\[
\pr(|X|>t)
=
\pr\!\left(e^{X^2/K_4^2}>e^{t^2/K_4^2}\right)
\le
e^{-t^2/K_4^2}\E e^{X^2/K_4^2}.
\]
By \emph{(iv)},
\[
\pr(|X|>t)\le 2e^{-t^2/K_4^2}.
\]
Thus \emph{(ii)} holds with $K_2\le K_4$.

\medskip
\noindent
\textbf{(iii)$\Rightarrow$(i).}
This is the step where centring matters. Let $X'$ be an independent copy of $X$.
Since $\E X=0$, Jensen's inequality gives
\[
e^{\lambda X}
=
e^{\lambda(X-\E X')}
\le
\E_{X'} e^{\lambda(X-X')}.
\]
Taking expectation in $X$,
\[
\E e^{\lambda X}\le \E e^{\lambda(X-X')}.
\]
Set
\[
Y:=X-X'.
\]
Then $Y$ is symmetric, so all odd moments of $Y$ vanish, and therefore
\[
\E e^{\lambda Y}
=
\sum_{m=0}^\infty \frac{\lambda^{2m}\E Y^{2m}}{(2m)!}.
\]
By Minkowski's inequality and \emph{(iii)},
\[
(\E |Y|^{2m})^{1/(2m)}
\le
(\E |X|^{2m})^{1/(2m)}+(\E |X'|^{2m})^{1/(2m)}
\le
2K_3\sqrt{2m}.
\]
Hence
\[
\E Y^{2m}\le (2K_3\sqrt{2m})^{2m}=(8K_3^2m)^m.
\]
Also,
\[
(2m)! = (1\cdots m)(m+1\cdots 2m)\ge m!\,m^m.
\]
Thus
\[
\frac{\lambda^{2m}\E Y^{2m}}{(2m)!}
\le
\frac{\lambda^{2m}(8K_3^2m)^m}{m!\,m^m}
=
\frac{(8K_3^2\lambda^2)^m}{m!}.
\]
Summing over $m$,
\[
\E e^{\lambda Y}
\le
\sum_{m=0}^\infty \frac{(8K_3^2\lambda^2)^m}{m!}
=
\exp(8K_3^2\lambda^2).
\]
Therefore
\[
\E e^{\lambda X}
\le
\exp(8K_3^2\lambda^2)
=
\exp\!\left(\frac{(4K_3)^2\lambda^2}{2}\right).
\]
So \emph{(i)} holds with
\[
K_1\le 4K_3.
\]

Combining the implications above shows that the four conditions are equivalent up to universal
constant factors.
\end{proof}

\begin{remark}
Sub-Gaussianity means that $X$ behaves, from the point of view of tails and moments, like a Gaussian random variable.  
The four formulations say the same thing in different languages: mgf, tails, moments, and Orlicz norms.
\end{remark}

We now state a one-sided version of the sub-exponential random variables.

\begin{de}[One-sided Bernstein condition / sub-exponential behavior]
A centered random variable $X$ satisfies the \emph{one-sided Bernstein condition} with parameters $(\sigma^2,c)$ if
\[
\E e^{\lambda X}
\le
\exp\!\left(\frac{\sigma^2\lambda^2}{2(1-c\lambda)}\right),
\qquad 0\le \lambda < \frac1c.
\]
When a corresponding two-sided version also holds for $-X$, one says that $X$ is \emph{sub-exponential}.
\end{de}

\begin{prop}[Equivalent characterizations of sub-exponentiality]
For a centered random variable $X$, the following are equivalent up to universal constants:
\begin{enumerate}[label=(\roman*)]
    \item $X$ satisfies a two-sided Bernstein mgf bound.
    \item There exists $K_1>0$ such that
    \[
    \pr(|X|>t)\le 2e^{-t/K_1},
    \qquad t\ge 0.
    \]
    \item There exists $K_2>0$ such that
    \[
    (\E |X|^p)^{1/p}\le K_2 p,
    \qquad p\ge 1.
    \]
    \item There exists $K_3>0$ such that
    \[
    \E e^{|X|/K_3}\le 2.
    \]
\end{enumerate}
\end{prop}

\begin{proof}[Proof sketch]
The proof parallels the sub-Gaussian case.

\textbf{(i)$\Rightarrow$(ii):} use Chernoff with the Bernstein mgf control.

\textbf{(ii)$\Rightarrow$(iii):} integrate the tail:
\[
\E |X|^p
=
\int_0^\infty p t^{p-1}\pr(|X|>t)\,dt
\le
2p\int_0^\infty t^{p-1}e^{-t/K_1}\,dt
=
2p! K_1^p.
\]
Thus $(\E |X|^p)^{1/p}\le C K_1 p$.

\textbf{(iii)$\Rightarrow$(iv):} expand $\E e^{|X|/L}$ into a power series and choose $L$ large.

\textbf{(iv)$\Rightarrow$(ii):} apply Markov:
\[
\pr(|X|>t)\le e^{-t/K_3}\E e^{|X|/K_3}\le 2e^{-t/K_3}.
\]
\end{proof}

\begin{cor}[Square of a sub-Gaussian variable is sub-exponential]
If $X$ is sub-Gaussian, then $X^2$ is sub-exponential.
\end{cor}

\begin{proof}
If $X$ is sub-Gaussian, then
\[
\E e^{X^2/K^2}\le 2
\]
for some $K>0$. This is exactly an exponential-integrability condition for $X^2$.
\end{proof}

The one-sided Bernstein condition results in a corresponding concentration inequality.
\begin{thm}[Bernstein upper-tail bound]
Let $X$ be a random variable satisfying the one-sided Bernstein condition with parameters $(\sigma^2,c)$:
\[
\E e^{\lambda(X-\E X)}
\le
\exp\!\left(\frac{\sigma^2\lambda^2}{2(1-c\lambda)}\right),
\qquad 0\le \lambda<\frac1c.
\]
Then for every $\varepsilon\ge 0$,
\[
\pr(X-\E X\ge \varepsilon)
\le
\exp\!\left(-\frac{\varepsilon^2}{2(\sigma^2+c\varepsilon)}\right).
\]
Equivalently, for every $\delta\in(0,1)$,
\[
\pr\!\left(
X-\E X
\ge
\sqrt{2\sigma^2\log(1/\delta)}+c\log(1/\delta)
\right)
\le \delta.
\]
\end{thm}

\begin{proof}
By Chernoff's bound, for any $0\le \lambda<1/c$,
\[
\pr(X-\E X\ge \varepsilon)
\le
e^{-\lambda\varepsilon}\E e^{\lambda(X-\E X)}
\le
\exp\!\left(
-\lambda\varepsilon+\frac{\sigma^2\lambda^2}{2(1-c\lambda)}
\right).
\]
Choose
\[
\lambda=\frac{\varepsilon}{\sigma^2+c\varepsilon}.
\]
Then $0\le \lambda<1/c$, and
\[
1-c\lambda = 1-\frac{c\varepsilon}{\sigma^2+c\varepsilon}
= \frac{\sigma^2}{\sigma^2+c\varepsilon}.
\]
Substituting gives
\[
-\lambda\varepsilon+\frac{\sigma^2\lambda^2}{2(1-c\lambda)}
=
-\frac{\varepsilon^2}{\sigma^2+c\varepsilon}
+\frac{\varepsilon^2}{2(\sigma^2+c\varepsilon)}
=
-\frac{\varepsilon^2}{2(\sigma^2+c\varepsilon)}.
\]
Hence
\[
\pr(X-\E X\ge \varepsilon)
\le
\exp\!\left(-\frac{\varepsilon^2}{2(\sigma^2+c\varepsilon)}\right).
\]

For the high-probability form, let $u=\log(1/\delta)$ and set
\[
\varepsilon=\sqrt{2\sigma^2u}+cu.
\]
Then
\[
\frac{\varepsilon^2}{2(\sigma^2+c\varepsilon)}\ge u,
\]
so the first inequality implies
\[
\pr(X-\E X\ge \varepsilon)\le e^{-u}=\delta.
\]
\end{proof}

\begin{remark}
Bernstein's inequality interpolates between two regimes:
\[
\varepsilon \ll \sigma^2/c
\quad \Rightarrow \quad
\exp\!\left(-\frac{\varepsilon^2}{2\sigma^2}\right)
\quad \text{(Gaussian regime),}
\]
and
\[
\varepsilon \gg \sigma^2/c
\quad \Rightarrow \quad
\exp\!\left(-\frac{\varepsilon}{2c}\right)
\quad \text{(exponential regime).}
\]
\end{remark}
Similar to the sub-exponential setting, the Bernstein condition is invariant under scaling and convex combinations.

\begin{prop}[Bernstein condition is preserved under nonnegative weighted sums]
Let $X_1,\dots,X_n$ be independent centered random variables, and suppose that each $X_i$ satisfies
\[
\E e^{\lambda X_i}
\le
\exp\!\left(\frac{\sigma_i^2\lambda^2}{2(1-c_i\lambda)}\right),
\qquad 0\le \lambda<1/c_i.
\]
Let $a_i\ge 0$ and define
\[
S=\sum_{i=1}^n a_i X_i.
\]
Then $S$ satisfies a one-sided Bernstein condition with parameters
\[
\sigma_S^2=\sum_{i=1}^n a_i^2\sigma_i^2,
\qquad
c_S=\max_{1\le i\le n} a_i c_i.
\]
\end{prop}

\begin{proof}
By independence,
\[
\E e^{\lambda S}
=
\prod_{i=1}^n \E e^{\lambda a_i X_i}
\le
\prod_{i=1}^n
\exp\!\left(\frac{a_i^2\sigma_i^2\lambda^2}{2(1-a_i c_i\lambda)}\right).
\]
If $0\le \lambda<1/c_S$, then $1-a_i c_i\lambda\ge 1-c_S\lambda$ for all $i$, so
\[
\E e^{\lambda S}
\le
\exp\!\left(
\sum_{i=1}^n
\frac{a_i^2\sigma_i^2\lambda^2}{2(1-c_S\lambda)}
\right)
=
\exp\!\left(
\frac{\sigma_S^2\lambda^2}{2(1-c_S\lambda)}
\right).
\]
\end{proof}

\begin{cor}[Bernstein inequality for sample means]
Let $X_1,\dots,X_n$ be i.i.d.\ with mean $\mu$, and assume that $X_1-\mu$ satisfies the one-sided Bernstein condition with parameters $(\sigma^2,c)$. Then
\[
\bar X_n-\mu
=
\frac1n\sum_{i=1}^n (X_i-\mu)
\]
satisfies the one-sided Bernstein condition with parameters
\[
\frac{\sigma^2}{n},
\qquad
\frac{c}{n}.
\]
Hence, for all $\varepsilon\ge 0$,
\[
\pr(\bar X_n-\mu\ge \varepsilon)
\le
\exp\!\left(
-\frac{n\varepsilon^2}{2(\sigma^2+c\varepsilon)}
\right).
\]
Equivalently, for every $\delta\in(0,1)$,
\[
\pr\!\left(
\bar X_n-\mu
\ge
\sqrt{\frac{2\sigma^2\log(1/\delta)}{n}}
+
\frac{c\log(1/\delta)}{n}
\right)\le \delta.
\]
\end{cor}

\begin{proof}
Apply the previous proposition with $a_i=1/n$. Then
\[
\sigma_S^2
=
n\cdot \frac{\sigma^2}{n^2}
=
\frac{\sigma^2}{n},
\qquad
c_S
=
\frac{c}{n}.
\]
Now apply Bernstein's upper-tail bound.
\end{proof}

\begin{cor}[Bounded observations: Bernstein and Hoeffding forms]
Let $X_1,\dots,X_n$ be i.i.d.\ with mean $\mu$.

\begin{enumerate}[label=(\alph*)]
    \item If $X_i-\mu\le c$ almost surely, then
    \[
    \pr(\bar X_n-\mu\ge \varepsilon)
    \le
    \exp\!\left(
    -\frac{n\varepsilon^2}{2(\var(X_1)+c\varepsilon)}
    \right).
    \]

    \item If $X_i\in[a,b]$ almost surely, then Hoeffding's inequality gives
    \[
    \pr(\bar X_n-\mu\ge \varepsilon)
    \le
    \exp\!\left(
    -\frac{2n\varepsilon^2}{(b-a)^2}
    \right).
    \]
\end{enumerate}
\end{cor}

\begin{proof}
Part (a) follows from the standard fact that bounded upper support implies a Bernstein-type mgf bound. Part (b) is Hoeffding's inequality.
\end{proof}

\begin{remark}
Hoeffding's inequality uses only boundedness. Bernstein's inequality uses both boundedness and variance, so it is often sharper when the variance is much smaller than the squared range.
\end{remark}

\section{Lecture 21: Concentration of other types}
\subsection{The bounded difference inequality}
We demonstrate that a function that is bounded in each of its coordinates also give rise to the concentration inequalities.
\begin{de}[Bounded differences property]
A function $f:\mathcal{X}^n\to\R$ satisfies the bounded differences property with constants $L_1,\dots,L_n$ if for every pair
$x=(x_1,\dots,x_n)$ and $x'=(x_1',\dots,x_n')$ that differ only in the $i$th coordinate,
\[
|f(x)-f(x')|\le L_i.
\]
\end{de}

\begin{thm}[McDiarmid's inequality]
Let $X_1,\dots,X_n$ be independent random variables and define
\[
Z=f(X_1,\dots,X_n),
\]
where $f$ satisfies bounded differences with constants $L_1,\dots,L_n$. Then for all $t\ge 0$,
\[
\pr(Z-\E Z\ge t)
\le
\exp\!\left(
-\frac{2t^2}{\sum_{i=1}^n L_i^2}
\right),
\]
and therefore
\[
\pr(|Z-\E Z|\ge t)
\le
2\exp\!\left(
-\frac{2t^2}{\sum_{i=1}^n L_i^2}
\right).
\]
Equivalently, $Z-\E Z$ is sub-Gaussian with variance proxy
\[
\frac14 \sum_{i=1}^n L_i^2.
\]
\end{thm}

\begin{proof}[Proof sketch]
Define the Doob martingale
\[
M_i=\E[Z\mid X_1,\dots,X_i],\qquad i=0,\dots,n,
\]
with $M_0=\E Z$ and $M_n=Z$. Then
\[
Z-\E Z = \sum_{i=1}^n (M_i-M_{i-1}).
\]
Changing only $X_i$ can change $f$ by at most $L_i$, so each martingale increment is almost surely bounded by $L_i$ in absolute value. Applying the Azuma--Hoeffding inequality to the martingale gives
\[
\pr(Z-\E Z\ge t)
\le
\exp\!\left(-\frac{2t^2}{\sum_i L_i^2}\right).
\]
Applying the same argument to $-Z$ yields the two-sided version.
\end{proof}

\begin{remark}
If no single observation can change the statistic very much, then the statistic is sharply concentrated around its mean.
\end{remark}


\begin{eg}[Kernel density estimation in $L^1$]
Let $X_1,\dots,X_n$ be i.i.d., let $K\ge 0$ be a kernel with $\int K(u)\,du=1$, and define
\[
\hat f_n(x)=\frac{1}{nh}\sum_{i=1}^n K\!\left(\frac{x-X_i}{h}\right).
\]
Set
\[
Z=\int_{\R} \bigl|\hat f_n(x)-\E \hat f_n(x)\bigr|\,dx.
\]
Then
\[
\pr\bigl(|Z-\E Z|\ge t\bigr)\le 2e^{-nt^2/2}.
\]
\end{eg}

\begin{proof}
If only $X_i$ is replaced by $X_i'$, then the estimator changes by
\[
\Delta_i(x)
=
\frac{1}{nh}
\left[
K\!\left(\frac{x-X_i}{h}\right)-K\!\left(\frac{x-X_i'}{h}\right)
\right].
\]
Hence
\[
\int |\Delta_i(x)|\,dx
\le
\frac{1}{nh}\int K\!\left(\frac{x-X_i}{h}\right)\,dx
+
\frac{1}{nh}\int K\!\left(\frac{x-X_i'}{h}\right)\,dx
=
\frac{2}{n}.
\]
So $Z$ satisfies bounded differences with $L_i=2/n$. Therefore
\[
\sum_{i=1}^n L_i^2 = n\cdot \frac{4}{n^2}=\frac{4}{n}.
\]
Applying McDiarmid's inequality gives
\[
\pr(|Z-\E Z|\ge t)
\le
2\exp\!\left(-\frac{2t^2}{4/n}\right)
=
2e^{-nt^2/2}.
\]
\end{proof}

\begin{remark}
Although the KDE is a function of all data points, each observation contributes only order $1/n$ in $L^1$, so the global estimator concentrates.
\end{remark}

\begin{eg}[Uniform deviation of the empirical measure]
Let $\mathcal{A}$ be a collection of measurable sets and define
\[
\mathbb{P}_n(A)=\frac1n\sum_{i=1}^n \1_{\{X_i\in A\}},
\qquad
P(A)=\pr(X\in A).
\]
Consider
\[
Z=\sup_{A\in\mathcal{A}} |\mathbb{P}_n(A)-P(A)|.
\]
Then
\[
\pr(|Z-\E Z|\ge t)\le 2e^{-2nt^2},
\]
and consequently, with probability at least $1-\delta$,
\[
Z\le \E Z+\sqrt{\frac{\log(2/\delta)}{2n}}.
\]
\end{eg}

\begin{proof}
If only the $i$th sample point is changed, then for every fixed $A$,
\[
\left|
\frac{1}{n}\sum_{j=1}^n \1_{\{X_j\in A\}}
-
\frac{1}{n}\sum_{j=1}^n \1_{\{X_j'\in A\}}
\right|
\le \frac{1}{n}.
\]
Taking the supremum over $A\in\mathcal{A}$ preserves this bound, so $L_i=1/n$ for all $i$. Hence
\[
\sum_{i=1}^n L_i^2 = n\cdot \frac{1}{n^2}=\frac1n.
\]
McDiarmid gives
\[
\pr(|Z-\E Z|\ge t)
\le
2\exp\!\left(-\frac{2t^2}{1/n}\right)
=
2e^{-2nt^2}.
\]
Set the right-hand side equal to $\delta$ and solve for $t$ to obtain the stated high-probability bound.
\end{proof}

\begin{remark}
The oscillation of the empirical measure over a whole class $\mathcal A$ is still controlled because changing one observation can alter every empirical probability by at most $1/n$.
\end{remark}

These tools are the standard bridge from probability inequalities to non-asymptotic statistics. They naturally generalize to the following result:
\begin{prop}[Concentration of bounded $U$-statistics]
Let $g:\mathcal X\times \mathcal X\to \R$ be symmetric and bounded:
\[
\|g\|_\infty \le b.
\]
Given i.i.d.\ $X_1,\dots,X_n$, define the order-$2$ $U$-statistic
\[
U_n
=
\frac{1}{\binom{n}{2}}
\sum_{1\le i<j\le n} g(X_i,X_j).
\]
Then
\[
\E U_n = \E g(X_1,X_2),
\]
and for every $t\ge 0$,
\[
\pr\bigl(|U_n-\E U_n|\ge t\bigr)
\le
2\exp\!\left(-\frac{n t^2}{8b^2}\right).
\]
\end{prop}

\begin{proof}
Unbiasedness is immediate by symmetry:
\[
\E U_n
=
\frac{1}{\binom{n}{2}}
\sum_{1\le i<j\le n}\E g(X_i,X_j)
=
\E g(X_1,X_2).
\]

To prove concentration, write $U_n=f(X_1,\dots,X_n)$. Replace only $X_i$ by $X_i'$. Then only the $(n-1)$ pairs involving index $i$ can change, so
\[
|f(x)-f(x')|
\le
\frac{n-1}{\binom{n}{2}}\cdot 2b
=
\frac{4b}{n}.
\]
Thus McDiarmid's inequality applies with $L_i=4b/n$ for all $i$. Since
\[
\sum_{i=1}^n L_i^2
=
n\cdot \frac{16b^2}{n^2}
=
\frac{16b^2}{n},
\]
we obtain
\[
\pr(|U_n-\E U_n|\ge t)
\le
2\exp\!\left(
-\frac{2t^2}{16b^2/n}
\right)
=
2\exp\!\left(-\frac{nt^2}{8b^2}\right).
\]
\end{proof}


\begin{prop}[Bounded empirical-process suprema]
Let $\mathcal F$ be a class of measurable functions such that
\[
\sup_{f\in\mathcal F}\|f\|_\infty \le b.
\]
Given i.i.d.\ $X_1,\dots,X_n$, define
\[
Z
=
\|\mathbb P_n-P\|_{\mathcal F}
:=
\sup_{f\in\mathcal F}
\left|
\frac1n\sum_{i=1}^n f(X_i)-\E f(X_1)
\right|.
\]
Then for every $t\ge 0$,
\[
\pr(Z\ge \E Z+t)
\le
\exp\!\left(-\frac{nt^2}{2b^2}\right).
\]
Equivalently, with probability at least $1-\delta$,
\[
Z
\le
\E Z + b\sqrt{\frac{2\log(1/\delta)}{n}}.
\]
\end{prop}

\begin{proof}
Write $Z=f(X_1,\dots,X_n)$. If only $X_i$ is replaced by $X_i'$, then for each fixed $f\in\mathcal F$,
\[
\left|
\frac1n\sum_{j=1}^n f(X_j) - \frac1n\sum_{j=1}^n f(X_j')
\right|
\le
\frac1n |f(X_i)-f(X_i')|
\le
\frac{2b}{n}.
\]
Taking the supremum over $f\in\mathcal F$ preserves the same bound:
\[
|Z(x)-Z(x')|\le \frac{2b}{n}.
\]
Thus McDiarmid applies with $L_i=2b/n$. Therefore
\[
\pr(Z\ge \E Z+t)
\le
\exp\!\left(
-\frac{2t^2}{n(2b/n)^2}
\right)
=
\exp\!\left(-\frac{nt^2}{2b^2}\right).
\]
Solving $\exp(-nt^2/(2b^2))=\delta$ yields the stated high-probability form.
\end{proof}

\begin{remark}
Even though $Z$ is a supremum over a possibly huge class $\mathcal F$, one observation still changes every empirical average by at most $2b/n$, which is enough for concentration.
\end{remark}

\subsection{Concentration of Lipschitz functions}
We demonstrate that Gaussian variables demonstrate sub-Gaussianity under Lipschitz transformations, which is quite generaly and admits many interesting applications.

\begin{thm}[Gaussian log-Sobolev inequality]
Let $\gamma_d$ denote the standard Gaussian measure on $\mathbb{R}^d$,
\[
d\gamma_d(x)=(2\pi)^{-d/2}e^{-\|x\|_2^2/2}\,dx.
\]
Then for every smooth function $g:\mathbb{R}^d\to\mathbb{R}$ such that $e^g$ and $|\nabla g|^2e^g$
are integrable with respect to $\gamma_d$,
\[
\Ent_{\gamma_d}(e^g)
:=
\int e^g g\,d\gamma_d
-
\Bigl(\int e^g\,d\gamma_d\Bigr)
\log\Bigl(\int e^g\,d\gamma_d\Bigr)
\le
\frac12\int |\nabla g|^2e^g\,d\gamma_d.
\]
Equivalently, for every smooth positive function $u$,
\[
\Ent_{\gamma_d}(u)
\le
\frac12\int \frac{|\nabla u|^2}{u}\,d\gamma_d.
\]
\end{thm}

\begin{proof}
Let $(P_t)_{t\ge 0}$ be the Ornstein--Uhlenbeck semigroup, defined by
\[
P_t f(x)
=
\mathbb{E}\Bigl[f\bigl(e^{-t}x+\sqrt{1-e^{-2t}}\,Z\bigr)\Bigr],
\qquad Z\sim N(0,I_d).
\]
Its generator is
\[
Lf=\Delta f-x\cdot \nabla f,
\]
and $\gamma_d$ is invariant under $P_t$, i.e.
\[
\int P_t f\,d\gamma_d=\int f\,d\gamma_d.
\]
Moreover, $P_t f\to \int f\,d\gamma_d$ as $t\to\infty$ for sufficiently integrable $f$.

We first recall the Gaussian integration by parts identity:
for smooth $f,h$ with sufficient decay,
\[
\int f\,Lh\,d\gamma_d
=
-\int \nabla f\cdot \nabla h\,d\gamma_d.
\]
Indeed, writing $d\gamma_d(x)=\rho(x)\,dx$ with
\[
\rho(x)=(2\pi)^{-d/2}e^{-\|x\|_2^2/2},
\qquad \nabla \rho(x)=-x\rho(x),
\]
we compute
\[
\int f\,Lh\,d\gamma_d
=
\int f(\Delta h-x\cdot \nabla h)\rho\,dx
=
-\int \nabla(f\rho)\cdot \nabla h\,dx-\int f\,x\cdot \nabla h\,\rho\,dx,
\]
and the terms involving $x\cdot \nabla h$ cancel, leaving
\[
\int f\,Lh\,d\gamma_d
=
-\int \rho\,\nabla f\cdot \nabla h\,dx
=
-\int \nabla f\cdot \nabla h\,d\gamma_d.
\]

Now let $u>0$ be smooth, and define
\[
F(t):=\Ent_{\gamma_d}(P_tu).
\]
Since $P_t$ preserves the Gaussian integral,
\[
\int P_tu\,d\gamma_d=\int u\,d\gamma_d=:m,
\]
and therefore
\[
F(t)=\int P_tu\log(P_tu)\,d\gamma_d-m\log m.
\]
Differentiating in $t$,
\[
F'(t)
=
\int \partial_t(P_tu)\log(P_tu)\,d\gamma_d
+
\int \partial_t(P_tu)\,d\gamma_d.
\]
The second term vanishes because $\int P_tu\,d\gamma_d$ is constant in $t$, and
$\partial_t P_tu=LP_tu$, so
\[
F'(t)=\int LP_tu\,\log(P_tu)\,d\gamma_d.
\]
Applying the Gaussian integration by parts identity with
$f=\log(P_tu)$ and $h=P_tu$, we obtain
\[
F'(t)
=
-\int \nabla \log(P_tu)\cdot \nabla(P_tu)\,d\gamma_d
=
-\int \frac{|\nabla P_tu|^2}{P_tu}\,d\gamma_d.
\]
Hence $F'(t)\le 0$, and integrating from $0$ to $\infty$ gives
\[
\Ent_{\gamma_d}(u)
=
\int_0^\infty \int \frac{|\nabla P_tu|^2}{P_tu}\,d\gamma_d\,dt,
\]
since $P_tu\to m$ as $t\to\infty$ and constants have zero entropy.

It remains to bound the integrand. By differentiating the Mehler formula under the expectation,
\[
\nabla P_tu = e^{-t}P_t(\nabla u).
\]
Therefore,
\[
\frac{|\nabla P_tu|^2}{P_tu}
=
e^{-2t}\frac{|P_t(\nabla u)|^2}{P_tu}.
\]
By Cauchy--Schwarz under the probability kernel defining $P_t$,
\[
|P_t(\nabla u)|^2
\le
P_t\!\left(\frac{|\nabla u|^2}{u}\right)P_tu,
\]
and thus
\[
\frac{|\nabla P_tu|^2}{P_tu}
\le
e^{-2t}P_t\!\left(\frac{|\nabla u|^2}{u}\right).
\]
Integrating with respect to $\gamma_d$ and using invariance of $\gamma_d$,
\[
\int \frac{|\nabla P_tu|^2}{P_tu}\,d\gamma_d
\le
e^{-2t}\int P_t\!\left(\frac{|\nabla u|^2}{u}\right)\,d\gamma_d
=
e^{-2t}\int \frac{|\nabla u|^2}{u}\,d\gamma_d.
\]
Substituting into the entropy representation,
\[
\Ent_{\gamma_d}(u)
\le
\int_0^\infty e^{-2t}\,dt
\int \frac{|\nabla u|^2}{u}\,d\gamma_d
=
\frac12\int \frac{|\nabla u|^2}{u}\,d\gamma_d.
\]

Finally, taking $u=e^g$, we have
\[
\nabla u=e^g\nabla g,
\qquad
\frac{|\nabla u|^2}{u}=|\nabla g|^2e^g,
\]
so
\[
\Ent_{\gamma_d}(e^g)\le \frac12\int |\nabla g|^2e^g\,d\gamma_d.
\]
This proves the claim.
\end{proof}

\begin{cor}[Sub-Gaussian mgf bound for Lipschitz functions]
Let $X\sim N(0,I_d)$, and let $f:\mathbb{R}^d\to\mathbb{R}$ be $L$-Lipschitz. Then
\[
\mathbb{E}\exp\!\bigl(\lambda(f(X)-\mathbb{E}f(X))\bigr)
\le
\exp\!\left(\frac{\lambda^2L^2}{2}\right),
\qquad \lambda\in\mathbb{R}.
\]
\end{cor}

\begin{proof}
Assume first that $f$ is smooth. Define
\[
M(t):=\mathbb{E}e^{tf(X)},
\qquad
\psi(t):=\log M(t).
\]
Applying the Gaussian log-Sobolev inequality with $g=tf$, we get
\[
\Ent_{\gamma_d}(e^{tf})
\le
\frac{t^2}{2}\int |\nabla f|^2e^{tf}\,d\gamma_d
\le
\frac{t^2L^2}{2}\,\mathbb{E}e^{tf(X)}.
\]
On the other hand,
\[
\Ent_{\gamma_d}(e^{tf})
=
\mathbb{E}[tf(X)e^{tf(X)}]-M(t)\log M(t)
=
M(t)\bigl(t\psi'(t)-\psi(t)\bigr).
\]
Hence
\[
t\psi'(t)-\psi(t)\le \frac{t^2L^2}{2}.
\]
For $t\neq 0$,
\[
\left(\frac{\psi(t)}{t}\right)'
=
\frac{t\psi'(t)-\psi(t)}{t^2}
\le
\frac{L^2}{2}.
\]
Integrating from $0$ to $\lambda>0$,
\[
\frac{\psi(\lambda)}{\lambda}
-
\lim_{t\downarrow 0}\frac{\psi(t)}{t}
\le
\frac{\lambda L^2}{2}.
\]
Since
\[
\lim_{t\downarrow 0}\frac{\psi(t)}{t}
=
\psi'(0)
=
\mathbb{E}f(X),
\]
it follows that
\[
\psi(\lambda)\le \lambda \mathbb{E}f(X)+\frac{\lambda^2L^2}{2},
\]
or equivalently,
\[
\mathbb{E}e^{\lambda(f(X)-\mathbb{E}f(X))}
\le
e^{\lambda^2L^2/2},
\qquad \lambda\ge 0.
\]
Applying the same argument to $-f$ yields the same estimate for $\lambda<0$.

For a merely Lipschitz $f$, one may approximate $f$ by smooth mollifications $f_\varepsilon$,
apply the estimate to $f_\varepsilon$, and then pass to the limit using dominated convergence.
\end{proof}


\begin{thm}[Gaussian concentration for Lipschitz functions]
Let $X\sim N(0,I_m)$ and let $f:\R^m\to \R$ be $L$-Lipschitz. Then
\[
f(X)-\E f(X)
\]
is sub-Gaussian with variance proxy $L^2$. In particular,
\[
\pr\bigl(|f(X)-\E f(X)|\ge t\bigr)
\le
2\exp\!\left(-\frac{t^2}{2L^2}\right),
\qquad t\ge 0.
\]
\end{thm}

\begin{proof}[Proof sketch]
The log-Sobolev inequality gives
\[
\E \exp\!\bigl(\lambda(f(X)-\E f(X))\bigr)
\le
\exp\!\left(\frac{\lambda^2L^2}{2}\right),
\qquad \lambda\in\R.
\]
Then Chernoff's bound gives
\[
\pr(f(X)-\E f(X)\ge t)
\le
\inf_{\lambda>0}
\exp\!\left(-\lambda t+\frac{\lambda^2L^2}{2}\right).
\]
Optimizing at $\lambda=t/L^2$ yields
\[
\pr(f(X)-\E f(X)\ge t)\le e^{-t^2/(2L^2)}.
\]
Apply the same bound to $-f$ and combine.
\end{proof}

\begin{remark}
For Gaussian inputs, Lipschitz observables behave as if they were one-dimensional Gaussians. The only parameter that matters is the Lipschitz constant.
\end{remark}

This result is powerful and has the following important corollaries:

\begin{cor}[Concentration of the Euclidean norm]
Let $X\sim N(0,I_n)$ and define
\[
f(x)=\frac{\norm{x}_2}{\sqrt n}.
\]
Then $f$ is $1/\sqrt n$-Lipschitz, and hence for every $t\ge 0$,
\[
\pr\left(
\left|
\frac{\norm{X}_2}{\sqrt n}
-
\E \frac{\norm{X}_2}{\sqrt n}
\right|
\ge t
\right)
\le
2\exp\!\left(-\frac{nt^2}{2}\right).
\]
\end{cor}

\begin{proof}
By the reverse triangle inequality,
\[
\left|
\frac{\norm{x}_2}{\sqrt n}
-
\frac{\norm{y}_2}{\sqrt n}
\right|
\le
\frac{\norm{x-y}_2}{\sqrt n}.
\]
Hence $f$ is $1/\sqrt n$-Lipschitz. Apply the Gaussian concentration theorem.
\end{proof}

\begin{cor}[Upper-tail bound for $\chi_n^2$]
Let $Y=\norm{X}_2^2\sim \chi_n^2$. Then for every $t\ge 0$,
\[
\pr\!\left(
\frac{Y}{n}\ge 1+2t+t^2
\right)
\le
\exp\!\left(-\frac{nt^2}{2}\right).
\]
\end{cor}

\begin{proof}
Since $\E \norm{X}_2\le \sqrt{\E \norm{X}_2^2}=\sqrt n$, the one-sided version of the previous result yields
\[
\pr\left(\frac{\norm{X}_2}{\sqrt n}\ge 1+t\right)
\le
\exp\!\left(-\frac{nt^2}{2}\right).
\]
Squaring the event gives
\[
\left\{
\frac{\norm{X}_2}{\sqrt n}\ge 1+t
\right\}
=
\left\{
\frac{Y}{n}\ge (1+t)^2
\right\}
=
\left\{
\frac{Y}{n}\ge 1+2t+t^2
\right\}.
\]
\end{proof}

\begin{remark}
Concentration of $\sqrt{Y}$ is often easier than concentration of $Y$ itself, because the square root is Lipschitz when there is a large deviation to the mean.
\end{remark}

\begin{cor}[Concentration of Gaussian order statistics]
Let $X_1,\dots,X_n$ be i.i.d.\ $N(0,1)$, and let
\[
X_{(1)}\le X_{(2)}\le \cdots \le X_{(n)}
\]
be the order statistics. Then for each $k\in\{1,\dots,n\}$,
\[
\pr\bigl(|X_{(k)}-\E X_{(k)}|\ge t\bigr)
\le
2\exp\!\left(-\frac{t^2}{2}\right),
\qquad t\ge 0.
\]
\end{cor}

\begin{proof}
It is a standard fact that the sorting map is non-expansive in $\ell_2$: if $x_{(k)}$ and $y_{(k)}$ denote the $k$th order statistics of $x$ and $y$, then
\[
|x_{(k)}-y_{(k)}|\le \norm{x-y}_2.
\]
Thus the function $f_k(x)=x_{(k)}$ is $1$-Lipschitz. Apply the Gaussian concentration theorem.
\end{proof}

\begin{remark}[Main idea]
Although order statistics are highly nonlinear, each individual order statistic is still a Lipschitz functional of the sample.
\end{remark}


\begin{cor}[Concentration of singular values]
Let $W\in \R^{n\times d}$ have i.i.d.\ $N(0,1)$ entries, and let
\[
\sigma_1(W)\ge \sigma_2(W)\ge \cdots \ge \sigma_{\min(n,d)}(W)\ge 0
\]
be its singular values. Then for every $k$ and every $t\ge 0$,
\[
\pr\left(
\left|
\sigma_k\!\left(\frac{W}{\sqrt n}\right)
-
\E \sigma_k\!\left(\frac{W}{\sqrt n}\right)
\right|
\ge t
\right)
\le
2\exp\!\left(-\frac{nt^2}{2}\right).
\]
\end{cor}

\begin{proof}
By Weyl's inequality,
\[
|\sigma_k(A)-\sigma_k(B)|\le \norm{A-B}_{\op}\le \norm{A-B}_F.
\]
Hence the map $A\mapsto \sigma_k(A)$ is $1$-Lipschitz with respect to the Frobenius norm. Therefore
\[
A\mapsto \sigma_k(A/\sqrt n)
\]
is $1/\sqrt n$-Lipschitz. Since $W$ is a standard Gaussian vector in $\R^{nd}$ after vectorization, the Gaussian concentration theorem applies.
\end{proof}

\begin{remark}
This is one of the most useful ways to get concentration for random matrices: identify the matrix quantity of interest as a Lipschitz function of the entries.
\end{remark}

\begin{lem}[Whitened covariance error]
Let $X_1,\dots,X_n\in \R^d$ be i.i.d.\ $N(0,\Sigma)$, and let
\[
\widehat\Sigma=\frac1n\sum_{i=1}^n X_iX_i^\top.
\]
Write
\[
X_i=\Sigma^{1/2}W_i,
\qquad
W_i\sim N(0,I_d).
\]
If $W\in \R^{n\times d}$ has rows $W_i^\top$, then
\[
\Sigma^{-1/2}(\widehat\Sigma-\Sigma)\Sigma^{-1/2}
=
\frac1n W^\top W-I_d.
\]
Consequently,
\[
\norm{\Sigma^{-1/2}(\widehat\Sigma-\Sigma)\Sigma^{-1/2}}_{\op}
=
\sup_{v\in S^{d-1}}
\left|
\frac1n\sum_{i=1}^n \bigl((v^\top W_i)^2-1\bigr)
\right|.
\]
\end{lem}

\begin{proof}
Using $X_i=\Sigma^{1/2}W_i$,
\[
\widehat\Sigma
=
\frac1n\sum_{i=1}^n \Sigma^{1/2}W_iW_i^\top\Sigma^{1/2}
=
\Sigma^{1/2}\left(\frac1n W^\top W\right)\Sigma^{1/2}.
\]
Subtract $\Sigma=\Sigma^{1/2}I_d\Sigma^{1/2}$ and whiten:
\[
\Sigma^{-1/2}(\widehat\Sigma-\Sigma)\Sigma^{-1/2}
=
\frac1nW^\top W-I_d.
\]
For a symmetric matrix $A$,
\[
\norm{A}_{\op}
=
\sup_{v\in S^{d-1}} |v^\top A v|.
\]
Applying this with $A=\frac1n W^\top W-I_d$ gives
\[
v^\top A v
=
\frac1n\sum_{i=1}^n (v^\top W_i)^2 - 1.
\]
\end{proof}

\begin{remark}
For each fixed $v$, the random variable $(v^\top W_i)^2-1$ is sub-exponential, so the fixed-direction fluctuations are easy. The hard part is the supremum over all $v\in S^{d-1}$, which is handled by empirical-process methods or by an $\varepsilon$-net argument.
\end{remark}

\section{Lecture 22: Applications in regression}
\subsection{The Gaussian sequence model and Hard thresholding}

\begin{de}[Gaussian sequence model]
The Gaussian sequence model is defined by the observations
\[
Y_j=\mu_j+\varepsilon_j,
\qquad j=1,\dots,d,\qquad \varepsilon_j \overset{\mathrm{iid}}{\sim} N\!\left(0,\frac{\sigma^2}{n}\right).
\]
\end{de}

\begin{lem}[Bias-variance decomposition of the risk]
For any estimator $\widehat\mu=\widehat\mu(Y)\in \R^d$,
\[
\E \norm{\widehat\mu-\mu}_2^2
=
\norm{\E\widehat\mu-\mu}_2^2
+
\tr\!\bigl(\Cov(\widehat\mu)\bigr).
\]
\end{lem}

\begin{proof}
Write
\[
\widehat\mu-\mu
=
(\widehat\mu-\E\widehat\mu)+(\E\widehat\mu-\mu).
\]
Squaring and taking expectations,
\[
\E \norm{\widehat\mu-\mu}_2^2
=
\E \norm{\widehat\mu-\E\widehat\mu}_2^2
+
\norm{\E\widehat\mu-\mu}_2^2
+
2\E\left\langle \widehat\mu-\E\widehat\mu,\E\widehat\mu-\mu\right\rangle.
\]
The cross term is zero because $\E(\widehat\mu-\E\widehat\mu)=0$. Also,
\[
\E \norm{\widehat\mu-\E\widehat\mu}_2^2
=
\tr(\Cov(\widehat\mu)).
\]
\end{proof}

\begin{cor}[Risk of the naive estimator]
For the estimator $\widehat\mu=Y$,
\[
\E \norm{Y-\mu}_2^2
=
\sum_{j=1}^d \E \varepsilon_j^2
=
\frac{\sigma^2 d}{n}.
\]
\end{cor}

\begin{proof}
Since $Y-\mu=\varepsilon$ and the coordinates are independent with variance $\sigma^2/n$,
\[
\E \norm{Y-\mu}_2^2
=
\sum_{j=1}^d \var(\varepsilon_j)
=
d\frac{\sigma^2}{n}.
\]
\end{proof}


Assume now that $\mu$ is sparse, with support
\[
S_\mu:=\{j:\mu_j\neq 0\},
\qquad s:=|S_\mu|.
\]

\begin{prop}[Oracle estimator when the support is known]
If the support $S_\mu$ were known, the oracle estimator
\[
\widehat\mu^{\,\mathrm{or}}_j
=
Y_j \1_{\{j\in S_\mu\}}
\]
would satisfy
\[
\E \norm{\widehat\mu^{\,\mathrm{or}}-\mu}_2^2
=
\frac{\sigma^2 s}{n}.
\]
\end{prop}

\begin{proof}
If $j\in S_\mu$, then $\widehat\mu^{\,\mathrm{or}}_j-\mu_j=Y_j-\mu_j=\varepsilon_j$.
If $j\notin S_\mu$, then $\mu_j=0$ and $\widehat\mu^{\,\mathrm{or}}_j=0$, so the error is zero.
Hence
\[
\E \norm{\widehat\mu^{\,\mathrm{or}}-\mu}_2^2
=
\sum_{j\in S_\mu}\E \varepsilon_j^2
=
\frac{\sigma^2 s}{n}.
\]
\end{proof}

\begin{remark}
This is the ideal benchmark.
If we knew exactly which coordinates were nonzero, we would only estimate those $s$ coordinates and ignore the rest.
The risk would then scale like $s/n$ instead of $d/n$.
Since usually $s\ll d$, this can be a dramatic improvement.
The real task is therefore to imitate this oracle behavior without knowing the support.
\end{remark}
To estimate the support, we applying hard thresholding that uses the observed magnitude $|Y_j|$ as a proxy for whether $\mu_j$ is nonzero.

\begin{de}[Hard-thresholding estimator]
For $\tau\ge 0$, define
\[
\widehat\mu^\tau_j
=
Y_j \1_{\{|Y_j|\ge \tau\}},
\qquad j=1,\dots,d.
\]
\end{de}


A convenient event is
\[
E_1
:=
\left\{
\max_{1\le j\le d} |\varepsilon_j|
\le \tau_0
\right\},
\qquad
\tau_0
=
\sigma \sqrt{\frac{2\log(2d/\delta)}{n}}.
\]

\begin{remark}
The event $E_1$ says that \emph{all} noise coordinates are simultaneously bounded by $\tau_0$.
The logarithmic term appears because we are controlling the maximum over $d$ coordinates, not just one coordinate.
This is the usual multiple-comparison price.
\end{remark}

The following lemma gives a clean high-probability event on which the rest of the argument becomes deterministic.
Once we know that every noise coordinate is at most $\tau_0$, we can compare $Y_j$ and $\mu_j$ coordinate by coordinate.

\begin{lem}[Uniform noise control]
For the choice above,
\[
\pr(E_1)\ge 1-\delta.
\]
\end{lem}

\begin{proof}
By the Gaussian tail bound,
\[
\pr(|\varepsilon_j|>\tau_0)
\le
2\exp\!\left(-\frac{n\tau_0^2}{2\sigma^2}\right)
=
2\exp\!\left(-\log\frac{2d}{\delta}\right)
=
\frac{\delta}{d}.
\]
Apply the union bound:
\[
\pr(E_1^c)
=
\pr\left(\max_{1\le j\le d}|\varepsilon_j|>\tau_0\right)
\le
\sum_{j=1}^d \pr(|\varepsilon_j|>\tau_0)
\le
d\cdot \frac{\delta}{d}
=
\delta.
\]
\end{proof}


\begin{prop}[Coordinatewise oracle inequality for hard thresholding]
Let $\widehat\mu^{2\tau_0}$ be the hard-threshold estimator with threshold $2\tau_0$. On the event $E_1$,
\[
\norm{\widehat\mu^{2\tau_0}-\mu}_2^2
\le
9\sum_{j=1}^d \min\{\mu_j^2,\tau_0^2\}.
\]
Therefore, with probability at least $1-\delta$,
\[
\norm{\widehat\mu^{2\tau_0}-\mu}_2^2
\le
9\sum_{j=1}^d \min\{\mu_j^2,\tau_0^2\}.
\]
\end{prop}

\begin{proof}
We work on $E_1$, so that $|\varepsilon_j|\le \tau_0$ for all $j$. We consider three cases.

\medskip
\noindent
\textbf{Case 1:} $|\mu_j|\le \tau_0$.

Then
\[
|Y_j|
\le
|\mu_j|+|\varepsilon_j|
\le
\tau_0+\tau_0
=
2\tau_0.
\]
Hence $\widehat\mu^{2\tau_0}_j=0$, so
\[
(\widehat\mu^{2\tau_0}_j-\mu_j)^2
=
\mu_j^2
=
\min\{\mu_j^2,\tau_0^2\}
\le
9\min\{\mu_j^2,\tau_0^2\}.
\]

\medskip
\noindent
\textbf{Case 2:} $|\mu_j|>3\tau_0$.

Then
\[
|Y_j|
\ge
|\mu_j|-|\varepsilon_j|
>
3\tau_0-\tau_0
=
2\tau_0,
\]
so $\widehat\mu^{2\tau_0}_j=Y_j$. Therefore
\[
(\widehat\mu^{2\tau_0}_j-\mu_j)^2
=
\varepsilon_j^2
\le
\tau_0^2
\le
9\min\{\mu_j^2,\tau_0^2\}.
\]

\medskip
\noindent
\textbf{Case 3:} $\tau_0<|\mu_j|\le 3\tau_0$.

In this regime,
\[
\widehat\mu^{2\tau_0}_j-\mu_j
=
-\mu_j \1_{\{|Y_j|<2\tau_0\}}
+
\varepsilon_j \1_{\{|Y_j|\ge 2\tau_0\}}.
\]
Hence
\[
(\widehat\mu^{2\tau_0}_j-\mu_j)^2
\le
\mu_j^2 \vee \varepsilon_j^2
\le
(3\tau_0)^2
=
9\tau_0^2
=
9\min\{\mu_j^2,\tau_0^2\}.
\]

Summing the three bounds over $j=1,\dots,d$ proves the claim.
\end{proof}

\begin{remark}
The reason that we threshold at $2\tau_0$ is that on the event $E_1$, a pure-noise coordinate satisfies $|Y_j|=|\varepsilon_j|\le \tau_0$, so it stays below $2\tau_0$ and gets removed; on the other hand, a sufficiently strong signal with $|\mu_j|>3\tau_0$ must satisfy
\[
|Y_j|\ge |\mu_j|-|\varepsilon_j|>2\tau_0,
\]
so it survives thresholding. Thus $2\tau_0$ creates a safety gap between noise-only coordinates and clearly nonzero coordinates.
\end{remark}

\begin{remark}
The magnitude of the bound
\[
\sum_{j=1}^d \min\{\mu_j^2,\tau_0^2\}
\]
has the following interpretations:
\begin{itemize}
    \item if $\mu_j$ is very small, the contribution is $\mu_j^2$, meaning it is cheap to ignore that coordinate;
    \item if $\mu_j$ is large, the contribution is only $\tau_0^2$, meaning the error is essentially bounded by the noise level.
\end{itemize}
So the estimator behaves almost like an oracle: it discards weak coordinates and keeps strong ones.
\end{remark}

\begin{remark}[Main idea]
The estimator behaves like an oracle up to a logarithmic threshold. Small coordinates are discarded, large coordinates are kept, and only the middle region requires a rough bound.
\end{remark}

\begin{cor}[Sparse high-probability rate]
If $\|\mu\|_0=s$, then with probability at least $1-\delta$,
\[
\norm{\widehat\mu^{2\tau_0}-\mu}_2^2
\le
9s\tau_0^2
=
18\,\frac{\sigma^2 s}{n}\log\!\left(\frac{2d}{\delta}\right).
\]
\end{cor}

\begin{proof}
If $\|\mu\|_0=s$, then only $s$ coordinates are nonzero, so
\[
\sum_{j=1}^d \min\{\mu_j^2,\tau_0^2\}
\le
s\tau_0^2.
\]
Now apply the previous proposition.
\end{proof}

\begin{remark}
This is the standard sparse recovery rate in this model.
Compared with the oracle rate $\sigma^2 s/n$, the extra factor $\log d$ is the cost of not knowing the support in advance and having to search among $d$ possible coordinates.
\end{remark}


\begin{cor}[Exact support recovery under a beta-min condition]
Assume
\[
\min_{j\in S_\mu} |\mu_j| > 3\tau_0.
\]
Then on the event $E_1$,
\[
\supp(\widehat\mu^{2\tau_0})=\supp(\mu).
\]
In particular, this holds with probability at least $1-\delta$.
\end{cor}

\begin{proof}
If $j\in S_\mu$, then $|\mu_j|>3\tau_0$, so on $E_1$,
\[
|Y_j|
\ge
|\mu_j|-|\varepsilon_j|
>
3\tau_0-\tau_0
=
2\tau_0.
\]
Hence $\widehat\mu^{2\tau_0}_j=Y_j\neq 0$.

If $j\notin S_\mu$, then $\mu_j=0$, so on $E_1$,
\[
|Y_j|=|\varepsilon_j|\le \tau_0<2\tau_0.
\]
Hence $\widehat\mu^{2\tau_0}_j=0$.

Thus the estimated and true supports coincide.
\end{proof}

\begin{remark}
The condition
\[
\min_{j\in S_\mu} |\mu_j| > 3\tau_0
\]
is a \emph{beta-min condition}: every nonzero signal must be large enough to stand out from the noise.
Without such a separation condition, exact support recovery is impossible in general, because very small nonzero coefficients are statistically indistinguishable from zero.
\end{remark}

\subsection{Soft thresholding and hyperparameter tuning}
Instead of using hard-thresholds, we may also use soft-thresholds.
\begin{de}[Soft-threshold map]
For $\lambda\ge 0$, define
\[
\eta_\lambda(y)
=
\begin{cases}
y-\lambda, & y>\lambda,\\[4pt]
0, & |y|\le \lambda,\\[4pt]
y+\lambda, & y<-\lambda.
\end{cases}
\]
Coordinatewise, for $y=(y_1,\dots,y_d)$,
\[
\eta_\lambda(y)
=
\bigl(\eta_\lambda(y_1),\dots,\eta_\lambda(y_d)\bigr).
\]
\end{de}

\begin{prop}[Soft thresholding is the lasso solution]
For any $Y\in\R^d$ and $\lambda\ge 0$,
\[
\widehat\beta^\lambda
:=
\argmin_{\beta\in\R^d}
\left\{
\frac12 \norm{Y-\beta}_2^2 + \lambda \norm{\beta}_1
\right\}
\]
is given coordinatewise by
\[
\widehat\beta^\lambda_j=\eta_\lambda(Y_j),
\qquad j=1,\dots,d.
\]
\end{prop}

\begin{proof}
The objective is separable:
\[
\frac12 \norm{Y-\beta}_2^2+\lambda\norm{\beta}_1
=
\sum_{j=1}^d
\left(
\frac12 (Y_j-\beta_j)^2+\lambda |\beta_j|
\right).
\]
Thus it suffices to solve the one-dimensional problem
\[
\min_{b\in\R}
\left\{
\frac12 (y-b)^2+\lambda |b|
\right\}.
\]

If $b>0$, differentiate:
\[
-(y-b)+\lambda=0
\quad\Longrightarrow\quad
b=y-\lambda,
\]
which is admissible only when $y>\lambda$.

If $b<0$, differentiate:
\[
-(y-b)-\lambda=0
\quad\Longrightarrow\quad
b=y+\lambda,
\]
which is admissible only when $y<-\lambda$.

If $|y|\le \lambda$, then $b=0$ satisfies the subgradient optimality condition
\[
0\in -(y-0)+\lambda[-1,1].
\]
Hence the minimizer is exactly
\[
b=
\begin{cases}
y-\lambda, & y>\lambda,\\
0, & |y|\le \lambda,\\
y+\lambda, & y<-\lambda.
\end{cases}
\]
Applying this coordinatewise completes the proof.
\end{proof}


\begin{remark}
This proposition shows that soft thresholding is not just a heuristic denoising rule:
it is exactly the solution of an $\ell_1$-penalized least-squares problem.
So the lasso can be viewed as a principled optimization formulation of thresholding.
\end{remark}


Now we assume the mean of the noisy observations is identical:
\[
Y \sim N(\mu, I_n).
\]
For a weakly differentiable estimator $\hat\mu(Y)$, Stein's unbiased risk estimate is
\[
\widehat R
=
-n + \|Y-\hat\mu(Y)\|_2^2 + 2\,\diver \hat\mu(Y),
\]
where
\[
\diver \hat\mu(Y)=\sum_{i=1}^n \frac{\partial \hat\mu_i(Y)}{\partial Y_i}.
\]
Then
\[
\E_\mu[\widehat R]=\E_\mu\|\hat\mu(Y)-\mu\|_2^2.
\]
So $\widehat R$ is an unbiased estimator of the true risk. 
For the soft-threshold estimator $\hat\mu^\lambda_i=\eta_\lambda(Y_i)$, we have
\[
Y_i-\eta_\lambda(Y_i)
=
\begin{cases}
\lambda, & Y_i>\lambda,\\
Y_i, & |Y_i|\le \lambda,\\
-\lambda, & Y_i<-\lambda,
\end{cases}
\]
hence
\[
\bigl(Y_i-\eta_\lambda(Y_i)\bigr)^2=\min\{Y_i^2,\lambda^2\}.
\]
Also,
\[
\frac{\partial}{\partial Y_i}\eta_\lambda(Y_i)=\1\{|Y_i|>\lambda\}
\]
away from the threshold points. Therefore,
\[
\widehat R(\lambda)
=
-n+\sum_{i=1}^n \min\{Y_i^2,\lambda^2\}
+2\sum_{i=1}^n \1\{|Y_i|>\lambda\}.
\]
Equivalently,
\[
\widehat R(\lambda)
=
\sum_{i=1}^n
\left[
1+\min\{Y_i^2,\lambda^2\}
-2\1\{|Y_i|\le \lambda\}
\right].
\]
Thus, one can choose
\[
\hat\lambda
=
\argmin_{\lambda\ge 0}\widehat R(\lambda),
\]
and then use $\hat\mu^{\hat\lambda}$ as the data-driven estimator. If $\widehat R(\lambda)$ is uniformly close to the true risk $R(\lambda)$ over a range of $\lambda$, then the selected $\hat\lambda$ behaves approximately like the oracle choice
\[
\lambda^\star=\argmin_{\lambda\ge 0} R(\lambda).
\]
So the SURE-selected estimator is asymptotically as good as the best estimator in the class.

\subsection{Lasso in the regression model}
Now we move to the regression model
\[
Y = X\beta^\star + \varepsilon,
\]
where $X\in\R^{n\times d}$ and $\beta^\star\in\R^d$. The least squares estimator is
\[
\hat\beta^{\mathrm{OLS}}
=
\argmin_{\beta\in\R^d}\|Y-X\beta\|_2^2.
\]
If $\mathrm{rank}(X)=d$ and $d\le n$, then
\[
\hat\beta^{\mathrm{OLS}}=(X^\top X)^{-1}X^\top Y.
\]

When $d>n$ or $\mathrm{rank}(X)<d$, the coefficient vector need not be unique, but the fitted values $X\hat\beta$ are unique.  
Among all least-squares solutions, the minimum-norm solution is
\[
\hat\beta^{\mathrm{MN}} = X^+Y,
\]
where $X^+$ is the Moore--Penrose pseudoinverse.

\begin{remark}
Even when $\beta$ is not unique, the prediction $X\beta$ may still be uniquely determined by the convex objective.
\end{remark}

A general penalised estimator has the form
\[
\hat\beta
=
\argmin_{\beta\in\R^d}
\left\{
\|Y-X\beta\|_2^2 + \lambda\,\mathrm{pen}(\beta)
\right\}.
\]
Typical penalties are:
\[
\ell_0:\ \|\beta\|_0,
\qquad
\ell_1:\ \|\beta\|_1,
\qquad
\ell_2:\ \|\beta\|_2^2.
\]
\begin{remark}
    There is also a constrained form:
\[
\hat\beta_t
=
\argmin_{\beta\in\R^d}
\|Y-X\beta\|_2^2
\quad\text{subject to}\quad
\mathrm{pen}(\beta)\le t.
\]
For convex penalties, the penalized and constrained formulations are closely related.
\end{remark}


The lasso estimator is
\[
\hat\beta
=
\argmin_{\beta\in\R^d}
\left\{
\frac{1}{2}\|Y-X\beta\|_2^2 + \lambda\|\beta\|_1
\right\}.
\]
Geometrically, the $\ell_1$ ball has corners, and these corners encourage sparse solutions. Because $\|\beta\|_1$ is not differentiable at $0$, optimality is expressed by subgradients:
\[
X^\top(Y-X\hat\beta)=\lambda s,
\qquad
s\in\partial\|\hat\beta\|_1.
\]
Coordinatewise,
\[
s_j=
\begin{cases}
\sign(\hat\beta_j), & \hat\beta_j\neq 0,\\
[-1,1], & \hat\beta_j=0.
\end{cases}
\]

Hence:
\begin{itemize}[leftmargin=2em]
    \item if $\hat\beta_j\neq 0$, then $X_j^\top(Y-X\hat\beta)=\lambda\sign(\hat\beta_j)$;
    \item if $\hat\beta_j=0$, then $|X_j^\top(Y-X\hat\beta)|\le \lambda$.
\end{itemize}


\begin{prop}
Consider the lasso problem
\[
\hat\beta \in \arg\min_{\beta\in\mathbb{R}^p}
\left\{
\frac12\|y-X\beta\|_2^2+\lambda\|\beta\|_1
\right\},
\qquad \lambda>0.
\]
Then the following hold:
\begin{enumerate}
    \item the fitted value $X\hat\beta$ is unique;
    \item the coefficient vector $\hat\beta$ need not be unique in general;
    \item the subgradient vector $s\in\partial \|\hat\beta\|_1$ associated with an optimum is unique.
\end{enumerate}
\end{prop}

\begin{proof}
Write
\[
f(\beta)=\frac12\|y-X\beta\|_2^2+\lambda\|\beta\|_1.
\]

\medskip
\noindent
\textbf{Uniqueness of the fitted value.}
Suppose $\hat\beta^{(1)}$ and $\hat\beta^{(2)}$ are two minimizers. Since $f$ is convex, for any $t\in[0,1]$,
\[
f\bigl(t\hat\beta^{(1)}+(1-t)\hat\beta^{(2)}\bigr)
\le
t f(\hat\beta^{(1)})+(1-t)f(\hat\beta^{(2)}),
\]
and the right-hand side equals the minimum value of $f$.
Hence equality must hold. Subtracting the convex function $\lambda\|\beta\|_1$ from the inequality yields that
\[
\left\|y-X\bigl(t\hat\beta^{(1)}+(1-t)\hat\beta^{(2)}\bigr)\right\|_2^2\ge t\left\|y-X\hat\beta^{(1)}\right\|_2^2+(1-t)\left\|y-X\hat\beta^{(2)}\right\|_2^2
\]
The strict convexity of the map
\[
u\mapsto \frac12\|y-u\|_2^2
\]
then implies that
\[
X\hat\beta^{(1)} = X\hat\beta^{(2)}.
\]
Thus all minimizers have the same fitted value, so $X\hat\beta$ is unique.

\medskip
\noindent
\textbf{Uniqueness of the subgradient vector.}
By the optimality condition for the lasso, $\hat\beta$ is optimal if and only if
\[
0 \in -X^\top(y-X\hat\beta)+\lambda\,\partial\|\hat\beta\|_1.
\]
Equivalently, there exists $s\in\partial\|\hat\beta\|_1$ such that
\[
X^\top(y-X\hat\beta)=\lambda s.
\]
Therefore,
\[
s=\frac{1}{\lambda}X^\top \left(y-X\hat \beta\right)
\]
is also unique. Thus the subgradient vector associated with any optimum is unique. This completes the proof.
\end{proof}
\begin{remark}
    We can also show that for all selection sets $E$ where $X$ has full rank on $E$, it admits a unique $\hat{\beta}$, i.e., for nice selection sets, Lasso only selects but does not estimate. Let $E=\supp(\hat\beta)$ and $\mathrm{rank}(X_E)=|E|$. Then any lasso solution satisfies
\[
\hat\beta_{E^c}=0,
\qquad
X_E^\top(Y-X_E\hat\beta_E)=\lambda s_E.
\]
Therefore,
\[
\hat\beta_E
=
(X_E^\top X_E)^+\bigl(X_E^\top Y-\lambda s_E\bigr)+v,
\qquad
v\in\nullsp(X_E).
\]
Since $\mathrm{rank}(X_E)=|E|$ implies $v=0$, we know that $\hat\beta$ is unique.
\end{remark}


\begin{remark}
    A sufficient condition for uniqueness is that the columns of $X$ are in \emph{general position}.  
Informally, this means that no signed column lies in the affine span of too few other signed columns.  
When the entries of $X$ come from a continuous distribution, the general position condition holds almost surely.
\end{remark}

\begin{prop}[Sparse support of a lasso solution]
For any $\lambda>0$, there exists a lasso solution with at most $n$ nonzero coefficients.
\end{prop}

This reflects the fact that the lasso chooses a sparse representation even when $d\gg n$.

\subsection{Hyperparameter tuning in Lasso}

To decide the hyperparameter $\lambda$, we need to evaluate the performance of the estimator $\hat{\beta}$. There are several different notions of error for $\hat\beta$:
\begin{enumerate}[leftmargin=2em]
    \item \textbf{In-sample prediction risk}
    \[
    \E\left[\frac1n\|X\hat\beta-X\beta^\star\|_2^2\right].
    \]

    \item \textbf{Estimation risk}
    \[
    \E\|\hat\beta-\beta^\star\|_2^2.
    \]

    \item \textbf{Out-of-sample prediction risk}
    for a new covariate vector $X_0$,
    \[
    \E\bigl[(X_0^\top\hat\beta-X_0^\top\beta^\star)^2\bigr].
    \]
\end{enumerate}
Usually, the in-sample prediction error is the easiest to analyse, since it depends only on the observed design matrix. \begin{prop}
Consider the normalised lasso
\[
\hat\beta
=
\argmin_{\beta\in\R^d}
\left\{
\frac{1}{2n}\|Y-X\beta\|_2^2 + \lambda\|\beta\|_1
\right\},
\]
under the model
\[
Y=X\beta^\star+\varepsilon.
\]
Let
\[
\Delta=\hat\beta-\beta^\star,
\]
and define the event
\[
\mathcal E_\lambda
=
\left\{
\left\|\frac{1}{n}X^\top\varepsilon\right\|_\infty
\le \frac{\lambda}{2}
\right\}.
\]
Then on the event \(\mathcal E_\lambda\),
\[
\frac{1}{n}\|X\hat\beta-X\beta^\star\|_2^2
=
\frac{1}{n}\|X\Delta\|_2^2
\le
3\lambda\|\beta^\star\|_1.
\]
In particular, this is a slow-rate oracle inequality for the in-sample prediction risk.

Moreover, if the columns of \(X\) satisfy
\[
\frac{1}{n}\|X_j\|_2^2\le 1,
\qquad j=1,\dots,d,
\]
and the noise is sub-Gaussian with scale parameter \(\sigma\), then with high probability
\[
\left\|\frac{1}{n}X^\top\varepsilon\right\|_\infty
\lesssim
\sigma\sqrt{\frac{\log(d/\delta)}{n}}.
\]
Hence a natural choice is
\[
\lambda \asymp \sigma \sqrt{\frac{\log(d/\delta)}{n}},
\]
for which
\[
\frac{1}{n}\|X\hat\beta-X\beta^\star\|_2^2
\lesssim
\sigma\|\beta^\star\|_1\sqrt{\frac{\log d}{n}},
\]
up to universal constants.
\end{prop}

\begin{proof}
By optimality of \(\hat\beta\),
\[
\frac{1}{2n}\|Y-X\hat\beta\|_2^2+\lambda\|\hat\beta\|_1
\le
\frac{1}{2n}\|Y-X\beta^\star\|_2^2+\lambda\|\beta^\star\|_1.
\]
Substituting \(Y=X\beta^\star+\varepsilon\) and writing \(\Delta=\hat\beta-\beta^\star\), we obtain
\[
\frac{1}{2n}\|\varepsilon-X\Delta\|_2^2+\lambda\|\hat\beta\|_1
\le
\frac{1}{2n}\|\varepsilon\|_2^2+\lambda\|\beta^\star\|_1.
\]
Expanding the square and cancelling \(\|\varepsilon\|_2^2/(2n)\) gives
\[
\frac{1}{2n}\|X\Delta\|_2^2
\le
\frac{1}{n}\varepsilon^\top X\Delta
+\lambda\bigl(\|\beta^\star\|_1-\|\hat\beta\|_1\bigr).
\]
Using Hölder's inequality,
\[
\frac{1}{n}\varepsilon^\top X\Delta
\le
\left\|\frac{1}{n}X^\top\varepsilon\right\|_\infty \|\Delta\|_1.
\]
Therefore,
\[
\frac{1}{2n}\|X\Delta\|_2^2
\le
\left\|\frac{1}{n}X^\top\varepsilon\right\|_\infty \|\Delta\|_1
+\lambda\bigl(\|\beta^\star\|_1-\|\hat\beta\|_1\bigr).
\]
On the event \(\mathcal E_\lambda\),
\[
\frac{1}{2n}\|X\Delta\|_2^2
\le
\frac{\lambda}{2}\|\Delta\|_1
+\lambda\bigl(\|\beta^\star\|_1-\|\hat\beta\|_1\bigr).
\]
Now apply the triangle inequality,
\[
\|\Delta\|_1=\|\hat\beta-\beta^\star\|_1\le \|\hat\beta\|_1+\|\beta^\star\|_1.
\]
Substituting this into the previous bound yields
\[
\frac{1}{2n}\|X\Delta\|_2^2
\le
\frac{\lambda}{2}\bigl(\|\hat\beta\|_1+\|\beta^\star\|_1\bigr)
+\lambda\bigl(\|\beta^\star\|_1-\|\hat\beta\|_1\bigr).
\]
After collecting terms,
\[
\frac{1}{2n}\|X\Delta\|_2^2
\le
\frac{3\lambda}{2}\|\beta^\star\|_1-\frac{\lambda}{2}\|\hat\beta\|_1
\le
\frac{3\lambda}{2}\|\beta^\star\|_1.
\]
Multiplying by \(2\), we conclude that
\[
\frac{1}{n}\|X\Delta\|_2^2
\le
3\lambda\|\beta^\star\|_1.
\]
Since \(X\Delta=X\hat\beta-X\beta^\star\), this proves the desired slow-rate prediction bound.

For the choice of \(\lambda\), note that if \(\|X_j\|_2^2/n\le 1\) for all \(j\) and the noise is sub-Gaussian with scale \(\sigma\), then standard maximal inequalities imply that, with high probability,
\[
\left\|\frac{1}{n}X^\top\varepsilon\right\|_\infty
\lesssim
\sigma\sqrt{\frac{\log(d/\delta)}{n}}.
\]
Thus taking
\[
\lambda \asymp \sigma \sqrt{\frac{\log(d/\delta)}{n}}
\]
ensures that the event \(\mathcal E_\lambda\) holds with high probability, and hence
\[
\frac{1}{n}\|X\hat\beta-X\beta^\star\|_2^2
\lesssim
\sigma\|\beta^\star\|_1\sqrt{\frac{\log d}{n}}.
\]
\end{proof}

\begin{remark}
This slow-rate bound does not require sparsity of \(\beta^\star\) or a strong structural condition on the design matrix \(X\). Its conclusion is only for prediction:
\[
\frac{1}{n}\|X\hat\beta-X\beta^\star\|_2^2
\lesssim
\lambda\|\beta^\star\|_1.
\]
It does not by itself imply a corresponding bound on \(\|\hat\beta-\beta^\star\|_1\) or \(\|\hat\beta-\beta^\star\|_2\).
\end{remark}

To obtain a sharper bound, we need stronger assumptions.

\begin{prop}
Suppose the spare model is given by
\[
Y=X\beta^\star+\varepsilon,\qquad S^\star=\operatorname{supp}(\beta^\star),
\qquad |S^\star|=s.
\]
and let the lasso estimator be defined by
\[
\hat\beta \in \arg\min_{\beta\in\mathbb{R}^d}
\left\{
\frac{1}{2n}\|Y-X\beta\|_2^2+\lambda\|\beta\|_1
\right\}.
\]
Define the estimation error
\[
\Delta=\hat\beta-\beta^\star,
\]
and the event
\[
\mathcal E_\lambda
=
\left\{
\left\|\frac{1}{n}X^\top\varepsilon\right\|_\infty\le \frac{\lambda}{2}
\right\}.
\]
Assume moreover that the restricted eigenvalue condition holds: there exists \(\kappa>0\) such that
\[
\frac{1}{n}\|Xu\|_2^2 \ge \kappa \|u\|_2^2
\qquad \text{for all } u\in C_3(S^\star),
\]
where
\[
C_3(S^\star)
=
\left\{
u\in\mathbb{R}^d:
\|u_{(S^\star)^c}\|_1\le 3\|u_{S^\star}\|_1
\right\}.
\]

Then on the event \(\mathcal E_\lambda\), one has
\[
\Delta\in C_3(S^\star),
\]
and
\[
\|\Delta\|_2 \le \frac{3\sqrt{s}\lambda}{\kappa},
\qquad
\|\Delta\|_1 \le \frac{12 s\lambda}{\kappa},
\qquad
\frac{1}{n}\|X\Delta\|_2^2 \le \frac{9 s\lambda^2}{\kappa}.
\]
In particular, if
\[
\lambda \asymp \sigma \sqrt{\frac{\log d}{n}},
\]
then
\[
\frac{1}{n}\|X\hat\beta-X\beta^\star\|_2^2
\lesssim
\frac{\sigma^2 s\log d}{n},
\]
up to universal constants and the restricted eigenvalue factor.
\end{prop}

\begin{proof}
Since \(\hat\beta\) minimizes the lasso objective and \(\beta^\star\) is feasible for comparison,
\[
\frac{1}{2n}\|Y-X\hat\beta\|_2^2+\lambda\|\hat\beta\|_1
\le
\frac{1}{2n}\|Y-X\beta^\star\|_2^2+\lambda\|\beta^\star\|_1.
\]
Using \(Y=X\beta^\star+\varepsilon\) and \(\Delta=\hat\beta-\beta^\star\), this becomes
\[
\frac{1}{2n}\|\varepsilon-X\Delta\|_2^2+\lambda\|\hat\beta\|_1
\le
\frac{1}{2n}\|\varepsilon\|_2^2+\lambda\|\beta^\star\|_1,
\]
hence
\[
\frac{1}{2n}\|X\Delta\|_2^2
\le
\frac{1}{n}\varepsilon^\top X\Delta+\lambda\|\beta^\star\|_1-\lambda\|\hat\beta\|_1.
\]
On the event \(\mathcal E_\lambda\),
\[
\frac{1}{n}\varepsilon^\top X\Delta
\le
\left\|\frac{1}{n}X^\top\varepsilon\right\|_\infty \|\Delta\|_1
\le
\frac{\lambda}{2}\|\Delta\|_1.
\]
Therefore
\[
\frac{1}{2n}\|X\Delta\|_2^2
\le
\frac{\lambda}{2}\|\Delta\|_1+\lambda\|\beta^\star\|_1-\lambda\|\hat\beta\|_1.
\]
Because \(\beta^\star_{(S^\star)^c}=0\), we have
\[
\|\beta^\star\|_1-\|\hat\beta\|_1
=
\|\beta^\star_{S^\star}\|_1-\|\hat\beta_{S^\star}\|_1-\|\hat\beta_{(S^\star)^c}\|_1
\le
\|\Delta_{S^\star}\|_1-\|\Delta_{(S^\star)^c}\|_1.
\]
Substituting this into the previous inequality gives
\[
\frac{1}{2n}\|X\Delta\|_2^2
\le
\frac{\lambda}{2}\bigl(\|\Delta_{S^\star}\|_1+\|\Delta_{(S^\star)^c}\|_1\bigr)
+\lambda\|\Delta_{S^\star}\|_1
-\lambda\|\Delta_{(S^\star)^c}\|_1.
\]
After collecting terms,
\[
\frac{1}{2n}\|X\Delta\|_2^2
\le
\frac{3\lambda}{2}\|\Delta_{S^\star}\|_1
-
\frac{\lambda}{2}\|\Delta_{(S^\star)^c}\|_1.
\]
Rearranging yields
\[
\frac{1}{2n}\|X\Delta\|_2^2
+
\frac{\lambda}{2}\|\Delta_{(S^\star)^c}\|_1
\le
\frac{3\lambda}{2}\|\Delta_{S^\star}\|_1.
\]
In particular,
\[
\|\Delta_{(S^\star)^c}\|_1\le 3\|\Delta_{S^\star}\|_1,
\]
so
\[
\Delta\in C_3(S^\star).
\]

Now apply the restricted eigenvalue condition:
\[
\frac{1}{n}\|X\Delta\|_2^2\ge \kappa \|\Delta\|_2^2.
\]
On the other hand, from the basic inequality above,
\[
\frac{1}{n}\|X\Delta\|_2^2
\le
3\lambda \|\Delta_{S^\star}\|_1
\le
3\lambda \sqrt{s}\,\|\Delta\|_2.
\]
Combining these gives
\[
\kappa \|\Delta\|_2^2
\le
3\lambda \sqrt{s}\,\|\Delta\|_2,
\]
and therefore
\[
\|\Delta\|_2 \le \frac{3\sqrt{s}\lambda}{\kappa}.
\]
Since \(\Delta\in C_3(S^\star)\),
\[
\|\Delta\|_1
=
\|\Delta_{S^\star}\|_1+\|\Delta_{(S^\star)^c}\|_1
\le
4\|\Delta_{S^\star}\|_1
\le
4\sqrt{s}\,\|\Delta\|_2
\le
\frac{12 s\lambda}{\kappa}.
\]
Finally,
\[
\frac{1}{n}\|X\Delta\|_2^2
\le
3\lambda \sqrt{s}\,\|\Delta\|_2
\le
\frac{9 s\lambda^2}{\kappa}.
\]

If now \(\lambda\asymp \sigma \sqrt{\log d/n}\), then
\[
\frac{1}{n}\|X\hat\beta-X\beta^\star\|_2^2
=
\frac{1}{n}\|X\Delta\|_2^2
\lesssim
\frac{\sigma^2 s\log d}{n},
\]
up to universal constants and the restricted eigenvalue factor.
\end{proof}

\begin{remark}
    This is the standard fast-rate prediction bound for the lasso. Under sparsity and a restricted eigenvalue condition, the prediction error scales as
\[
\frac{\sigma^2 s\log d}{n},
\]
which is substantially sharper than the slower approximation-type bounds that hold without such structural assumptions.

\end{remark}


\begin{remark}
    Prediction consistency is weaker than exact variable-selection consistency. To recover the support $S^\star$ exactly, one usually needs:
\begin{itemize}[leftmargin=2em]
    \item a design condition stronger than RE, often an irrepresentable-type condition;
    \item a \emph{beta-min} condition ensuring that the nonzero coefficients are large enough:
    \[
    \min_{j\in S^\star} |\beta_j^\star|
    \quad\text{must dominate the estimation error level.}
    \]
\end{itemize}

Thus, exact support recovery is much harder than obtaining a small prediction error.
\end{remark}

\begin{remark}[Choosing the tuning parameter in practice]
    Two standard strategies were emphasized:
\begin{enumerate}[leftmargin=2em]
    \item \textbf{SURE}, when a Gaussian sequence-model approximation is appropriate.
    \item \textbf{Cross-validation}, especially in regression problems where prediction performance is the main criterion.
\end{enumerate}

The tuning parameter controls the amount of regularization:
\begin{itemize}[leftmargin=2em]
    \item too small $\lambda$ gives instability and overfitting;
    \item too large $\lambda$ gives excessive shrinkage and underfitting.
\end{itemize}
\end{remark}







More generally, if the true regression function is not exactly linear, say
\[
Y=f^\star + \varepsilon,
\]
and we fit the constrained lasso
\[
\hat\beta_t
=
\argmin_{\|\beta\|_1\le t}\frac1n\|Y-X\beta\|_2^2,
\]
then one can show an approximation-type bound of the form
\[
\frac1n\|X\hat\beta_t-f^\star\|_2^2
\lesssim
\inf_{\|\beta\|_1\le t}\frac1n\|X\beta-f^\star\|_2^2
+
\sigma t\sqrt{\frac{\log d}{n}}.
\]
So the lasso trades off approximation error and stochastic error.

\begin{remark}[Connection to nonparametric regression]
   A useful way to connect nonparametric regression with high-dimensional linear methods is through a basis expansion. Let $Z_i\in\R$ be a scalar input, and choose basis functions
\[
\phi_1,\phi_2,\dots,\phi_m.
\]
Define the feature vector
\[
X_i = \bigl(\phi_1(Z_i),\phi_2(Z_i),\dots,\phi_m(Z_i)\bigr)\in\R^m.
\]
Then any function represented in this basis has the form
\[
f_\beta(z)=\sum_{j=1}^m \beta_j \phi_j(z),
\]
and the regression model becomes
\[
Y_i = f_0(Z_i)+\varepsilon_i
\approx X_i^\top \beta + \varepsilon_i.
\]

This converts a function-estimation problem into a linear regression problem in a possibly high-dimensional dictionary.  
If the function is well approximated by only a few basis coefficients, then sparse methods such as the Lasso are natural. 
\end{remark}

However, Lasso can still give an oracle inequality in this case.

\begin{prop}
Assume the data obey
\[
Y=f^\star+\varepsilon \in \mathbb{R}^n,
\]
where \(f^\star\in\mathbb{R}^n\) is an arbitrary regression signal and
\(\varepsilon\in\mathbb{R}^n\) is a noise vector. Let
\[
\hat\beta_t \in \arg\min_{\|\beta\|_1\le t}\frac{1}{n}\|Y-X\beta\|_2^2,
\]
where \(X\in\mathbb{R}^{n\times d}\) is the design matrix.

Then for every \(t\ge 0\), with probability at least \(1-c_1 d^{-c_2}\),
\[
\frac{1}{n}\|X\hat\beta_t-f^\star\|_2^2
\le
\inf_{\|\beta\|_1\le t}\frac{1}{n}\|X\beta-f^\star\|_2^2
+
C\,\sigma t\sqrt{\frac{\log d}{n}},
\]
provided that the columns of \(X\) are normalized so that
\[
\frac{1}{n}\|X_j\|_2^2\le 1,\qquad j=1,\dots,d,
\]
and the noise coordinates are mean-zero sub-Gaussian with scale parameter \(\sigma\).
\end{prop}

\begin{proof}
Since \(\hat\beta_t\) minimizes the constrained empirical risk over the set
\(\{\beta:\|\beta\|_1\le t\}\), for every \(\beta\) with \(\|\beta\|_1\le t\),
\[
\frac{1}{n}\|Y-X\hat\beta_t\|_2^2
\le
\frac{1}{n}\|Y-X\beta\|_2^2.
\]
Substituting \(Y=f^\star+\varepsilon\), we get
\[
\frac{1}{n}\|f^\star+\varepsilon-X\hat\beta_t\|_2^2
\le
\frac{1}{n}\|f^\star+\varepsilon-X\beta\|_2^2.
\]
Expanding both sides and cancelling the common term \(\|\varepsilon\|_2^2/n\) yields
\[
\frac{1}{n}\|X\hat\beta_t-f^\star\|_2^2
\le
\frac{1}{n}\|X\beta-f^\star\|_2^2
+
\frac{2}{n}\varepsilon^\top X(\hat\beta_t-\beta).
\]
Using Hölder's inequality,
\[
\frac{2}{n}\varepsilon^\top X(\hat\beta_t-\beta)
\le
2\left\|\frac{1}{n}X^\top\varepsilon\right\|_\infty
\|\hat\beta_t-\beta\|_1.
\]
Since both \(\hat\beta_t\) and \(\beta\) belong to the \(\ell_1\)-ball of radius \(t\),
\[
\|\hat\beta_t-\beta\|_1
\le
\|\hat\beta_t\|_1+\|\beta\|_1
\le
2t.
\]
Therefore,
\[
\frac{1}{n}\|X\hat\beta_t-f^\star\|_2^2
\le
\frac{1}{n}\|X\beta-f^\star\|_2^2
+
4t\left\|\frac{1}{n}X^\top\varepsilon\right\|_\infty.
\]
Taking the infimum over all \(\beta\) with \(\|\beta\|_1\le t\), we obtain the deterministic bound
\[
\frac{1}{n}\|X\hat\beta_t-f^\star\|_2^2
\le
\inf_{\|\beta\|_1\le t}
\frac{1}{n}\|X\beta-f^\star\|_2^2
+
4t\left\|\frac{1}{n}X^\top\varepsilon\right\|_\infty.
\]
It remains to control the stochastic term. Under the column normalization, \(\|X_j\|_2^2/n\le 1\) and the sub-Gaussian assumption on \(\varepsilon\), a standard maximal inequality gives
\[
\left\|\frac{1}{n}X^\top\varepsilon\right\|_\infty
\le
C'\sigma\sqrt{\frac{\log d}{n}}
\]
with probability at least \(1-c_1 d^{-c_2}\). Substituting this estimate into the previous display gives
\[
\frac{1}{n}\|X\hat\beta_t-f^\star\|_2^2
\le
\inf_{\|\beta\|_1\le t}
\frac{1}{n}\|X\beta-f^\star\|_2^2
+
C\,\sigma t\sqrt{\frac{\log d}{n}},
\]
for a universal constant \(C>0\). This proves the claim.
\end{proof}

\begin{remark}
The bound shows that the constrained lasso balances two effects:
\begin{itemize}
    \item the \emph{approximation error}
    \[
    \inf_{\|\beta\|_1\le t}\frac{1}{n}\|X\beta-f^\star\|_2^2,
    \]
    which measures how well the true regression function can be approximated by a linear predictor with \(\ell_1\)-norm at most \(t\);
    \item the \emph{stochastic error}
    \[
    \sigma t\sqrt{\frac{\log d}{n}},
    \]
    which grows with the size of the feasible set.
\end{itemize}
Hence larger values of \(t\) reduce approximation error but increase stochastic error, while smaller values of \(t\) do the reverse.
\end{remark}


